%% File: semicirculos.Rnw
%% R/exams template for the semicircles problem

\extype{schoice}         % tipo: opción múltiple con una sola solución
\exname{Semicírculos}    % nombre interno del ejercicio
\exsolution{C}           % la respuesta correcta (índice de la lista 1=A,2=B,…)
\exshuffle{TRUE}         % barajar las opciones

\begin{question}
En la figura adjunta, $AC=12\text{ cm}$ y $AO = 2\,BP$. Se han dibujado dos semicírculos de centro en $O$ y $B$ sobre el segmento $AC$, tal como se muestra.

¿Cuáles son las medidas de los radios de los dos semicírculos?

\begin{center}
\begin{tikzpicture}[scale=0.5]
  % -------------------
  % Primera gráfica
  % -------------------
  \begin{scope}[yshift=4cm]
    % puntos
    \coordinate (A) at (0,0);
    \coordinate (O) at (4,0);
    \coordinate (B) at (8,0);
    \coordinate (P) at (10,0);
    \coordinate (C) at (12,0);
    % línea base
    \draw (A) -- (C);
    % puntos etiquetados
    \foreach \pt/\name in {A/A,O/O,B/B,P/P,C/C} {
      \fill (\pt) circle (1.5pt) node[below] {\name};
    }
    % semicírculos
    \draw (O) ++(4,0) arc[start angle=0,end angle=180,radius=4];
    \draw (B) ++(2,0) arc[start angle=0,end angle=180,radius=2];
    % etiquetas de longitudes
    \node[above] at ($(A)!0.5!(C)$) {$AC=12\text{ cm}$};
    \node[above] at ($(A)!0.5!(O)$) {$AO$};
    \node[above] at ($(B)!0.5!(P)$) {$BP$};
  \end{scope}

  % -------------------
  % Segunda gráfica (esquema con X)
  % -------------------
  \begin{scope}
    % puntos
    \coordinate (A) at (0,0);
    \coordinate (O) at (2,0);
    \coordinate (B) at (4,0);
    \coordinate (P) at (5,0);
    \coordinate (C) at (6,0);
    % líneas discontinuas verticales y superior
    \draw[dashed] (A) -- ++(0,2.3) -- ++(6,0) -- ++(0,-2.3);
    \draw[dashed] (O) -- ++(0,1.8);
    \draw[dashed] (B) -- ++(0,1.3);
    \draw[dashed] (P) -- ++(0,1.3);
    % línea base con etiqueta total
    \draw (A) -- (C) node[midway,above] {12 cm};
    % puntos y etiquetas
    \foreach \pt/\name in {A/A,O/O,B/B,P/P,C/C} {
      \fill (\pt) circle (1.5pt) node[below] {\name};
    }
    % semicírculos
    \draw (O) ++(2,0) arc[start angle=0,end angle=180,radius=2];
    \draw (B) ++(1,0) arc[start angle=0,end angle=180,radius=1];
    % etiquetas de los segmentos
    \draw[<->] (A) -- (O) node[midway,below] {$2X$ cm};
    \draw[<->] (O) -- (B) node[midway,below] {$2X$ cm};
    \draw[<->] (B) -- (P) node[midway,below] {$X$ cm};
    \draw[<->] (P) -- (C) node[midway,below] {$X$ cm};
  \end{scope}

\end{tikzpicture}
\end{center}

\begin{enumerate}[A)]
\item $2\text{ cm}$ y $2\text{ cm}$, respectivamente.
\item $2\text{ cm}$ y $4\text{ cm}$, respectivamente.
\item $4\text{ cm}$ y $2\text{ cm}$, respectivamente.
\item $4\text{ cm}$ y $4\text{ cm}$, respectivamente.
\end{enumerate}
\end{question}

\begin{solution}
Se tiene la partición
\[
AO + OB + BP + PC = 12,\quad
AO=2X,\;OB=2X,\;BP=X,\;PC=X.
\]
Por tanto $2X+2X+X+X=6X=12\implies X=2$. El radio del semicírculo de centro $O$ es
\[
OB=2X=4\text{ cm},
\]
y el radio del semicírculo de centro $B$ es
\[
BP=X=2\text{ cm}.
\]
\end{solution}

\begin{explanations}
Detallamos los pasos:
1) Etiquetamos los cuatro trozos de la recta como $2X,2X,X,X$;
2) Al sumar deben dar $12$ cm, de donde $6X=12\implies X=2$;
3) Así $OB=2X=4$ cm y $BP=X=2$ cm.
\end{explanations}
